\documentclass[11pt, a4paper]{article} 
\usepackage[margin = 1in]{geometry} 
\usepackage[]{amsmath, amssymb, amsfonts, amsthm} 
\allowdisplaybreaks
\usepackage[shortlabels]{enumitem} 
\usepackage[]{mathtools} 

\NewDocumentCommand{\htb}{}{\hat{\beta}}
\NewDocumentCommand{\Var}{}{\operatorname{Var}}
\NewDocumentCommand{\rank}{}{\operatorname{rank}}
\NewDocumentCommand{\ginv}{m}{#1^{-}}

\parindent 0pt

\title{Simple Linear Model}
\author{Notes by Moutushi Chatterji \LaTeX ed by Sathvik H S}
\date{}

\begin{document}
\maketitle

\section*{Model Assumptions and Model Objectives}

Consider the general linear model of the form 
\begin{equation}
	Y = X \beta + \varepsilon \label{model_eq}
\end{equation}
where 
\[ Y = 
\begin{pmatrix}
    y_{1} \\ 
    y_{2} \\ 
    \vdots \\ 
    y_{n}
\end{pmatrix}
\quad
X = 
\begin{pmatrix}
x_{11} & x_{12} & \cdots & x_{1p} \\
x_{21} & x_{22} & \cdots & x_{2p} \\
\vdots & \vdots & \ddots & \vdots \\
x_{n1} & x_{n2} & \cdots & x_{np} 
\end{pmatrix}
\quad
\beta =
\begin{pmatrix}
\beta_{1} \\
\beta_{2} \\
\vdots \\
\beta_{p}
\end{pmatrix}
\quad
\varepsilon = 
\begin{pmatrix}
    \varepsilon_{1} \\ 
    \varepsilon_{2} \\ 
    \vdots \\ 
    \varepsilon_{n}
\end{pmatrix}
\] 

Here \( Y \) is the vector of \( n \) observations, \( \beta \) is the vector of parameters, \( \varepsilon \) is the vector of random errors and \( X \) is the design matrix.
Both \( X \) and \( \beta \) are fixed.
\subsection*{Assumptions}
We assume \( \varepsilon \) to have the following properties:
\begin{align}
	\mathbb{E}(\varepsilon) &= 0 \\
	\Var(\varepsilon) &= \sigma^{2} I_{n}
\end{align}
The second assumption implies that \( \varepsilon_{1}, \varepsilon_{2}, \dots, \varepsilon_{n} \) have teh same but unknown variances \( \sigma^{2} \) and are uncorrelated.

\subsection*{Objectives}
\begin{enumerate}[i)]
	\item To estimate the unknown parameters \( \beta_{1}, \beta_{2}, \dots \beta_{p} \) is possible OR at least estimate those linear combinations of \( \beta \) that can be estimated.
	\item To test suitable statistical hypotheses about \( \beta \) OR at least functions of \( \beta \).
\end{enumerate}

\section*{Derivation}
The rank of \( X \) is assumed to \( r \) such that \( r \leq \min (n, p) \). If \( r = p \leq n \), the model \eqref{model_eq} is said to be a \textit{full rank model}, otherwise it is called a non full rank model.

Define \( \htb \) as a function of known quantities like \( Y, X, \) etc. such that \( \htb \) is ``close'' to \( \beta \) in some sense. Then, 
\begin{enumerate}[1.]
	\item \( e = Y - X\htb \) is the \textit{vector of residuals}.
	\item \( \varepsilon = Y - X\beta \) is the \textit{vector of errors}.
\end{enumerate}

One way to obtain \( \htb \) is to use \textit{method of least squares}.

Consider the sum of squares of residuals, i.e, 
\begin{align*}
	e'e &= (Y - X\htb)'(Y - X\htb) \\
	    &= Y'Y - Y'X\htb - \htb'X'Y - \htb'X'X\htb \\
	    &= Y'Y - 2\htb X'Y + \htb'X'X\htb
\end{align*}
The last equality follows since \( Y'X\htb = \htb'X'Y \) as both sides are \( 1 \times 1 \) matrices.
\[ \implies \frac{\partial }{\partial \htb} (e'e) = -2X'Y + 2(X'X)\htb \] 
Equating this to \( 0 \) we have, 
\begin{equation}
	X'Y = (X'X) \htb \label{normal_eq}
\end{equation}
The above collection of equations are called \textit{normal equations}.

When the inverse of \( (X'X) \) exists, we simply have
\begin{equation}
	\htb = (X'X)^{-1}X'Y
\end{equation}

\subsection*{Some Results}

\begin{enumerate}[1.]
	\item The matrix \( S = X'X \) is a symmetric matrix having same rank as \( X \).
\begin{proof}
If a vector \( \alpha \) is orthogonal to the rows of the matrix \( X \), then \( X \alpha = 0 \implies X'X \alpha = 0 \), i.e, \( \alpha \) is orthogonal to the rows of \( X'X \) as well.

Conversely, if \( X'X \alpha = 0 \), then \( \alpha'X'X\alpha = 0\) or \( Z'Z = 0 \) where \( Z = X\alpha \). But \( Z'Z = z_{1}^{2} + z_{2}^{2} + \dots + z_{n}^{2} \) which implies \( Z = X\alpha = 0 \), i.e, \( \alpha \) is orthogonal to the rows of \( X \).

Thus \( X \) and \( X'X \) has the same rank \( r \). This also shows that the row space of \( X \) and row space of \( S = X'X \) are the same.
\end{proof}

	\item If \( \htb \) be any particular solution to the \eqref{normal_eq}, then it minimizes \( e'e \), the sum of square of residuals.
\begin{proof}
Consider any other values of \( \beta \), say \( \beta_{0} \). Then,
\begin{align*}
	(Y-X\beta_{0})'(Y-X\beta_{0}) &= (Y - X\htb + X\htb - X\beta_{0})'(Y - X\htb + X\htb - X\beta_{0}) \\
				      &= (Y - X\htb)'(Y - X\htb) + \left\{ X(\htb - \beta_{0}) \right\} ' (Y-X\htb) \\
				      &\quad + (Y-X\htb)' \left\{ X(\htb - \beta_{0}) \right\} + \left\{ X(\htb - \beta_{0}) \right\}' \left\{ X(\htb - \beta_{0}) \right\} \\
				      &= (Y - X\htb)'(Y - X\htb) + (\htb - \beta_{0})'X'(Y-X\htb) \\
				      &\quad + (Y-X\htb)'X(\htb-\beta_{0}) + \left\{ X(\htb - \beta_{0}) \right\}' \left\{ X(\htb - \beta_{0}) \right\} \\
				      &= (Y - X\htb)'(Y - X\htb) + (\htb - \beta_{0})'(X'Y - X'X\htb) \\
				      &\quad + (X'Y - X'X\htb)'(\htb-\beta_{0}) + m'm \tag*{where \( m = X(\htb - \beta_{0}) \)} \\
				      &= (Y - X\htb)'(Y - X\htb) + m'm \tag*{(from \eqref{normal_eq})} \\
				      & \geq (Y - X\htb)'(Y - X\htb) \qedhere
\end{align*}
\end{proof}

Thus the minimized value can be unambiguously denoted by \( SSE = (Y - X\htb)'(Y - X\htb) \) and stands for \textit{sum of squares due to error}.
\end{enumerate}

\section*{Generalized Inverse and It's Properties}
\subsection*{Definition}
\begin{enumerate}[1.]
	\item Any \( n \times m \) matrix \( \ginv{A} \) satisying the relation \( A\ginv{A}A = A \) is defined as a \textit{generalized inverse (g-inverse)} of the \( m \times n \) matrix \( A \).
	\item An \( n \times m \) matrix \( \ginv{A} \) is defined to be a \textit{generalized inverse} of the \( m \times n \) matrix \( A \) if for every \( u \) satisfying \( \rank(A) = \rank(A|u) \), \( \ginv{A}u \) is a solution of the equations \( Ax = u \) for unknown \( u \). 
		In other words, \( \ginv{A}u \) is a solution of the equations \( Ax = u \) for unknown \( u \) whenever the system is \textit{consistent}.
\end{enumerate}

\subsection*{Properties}
Define an \( n \times n \) matrix \( H = \ginv{A}A \). Then, 
\begin{enumerate}[1.]
	\item \( AH = A \)
	\item \( H^{2} = H \)
	\item \( \operatorname{Tr}(H) = \rank(H) = \rank(A) \)
\end{enumerate}
\begin{proof}
The first two properties follow directly from definition. 

The first equality in property 3 follows as \( H \) is an idempotent matrix (property 2) and is similar to a diagonal matrix containing eigenvalues of \( H \) (spectral theorem). But the only eigenvalues of an idempotent matrix are \( 0 \) and \( 1 \) as the polynomial \( x^2 - x \) annihilates \( H \).

For the second equality, observe that \( \ginv{A} A \alpha = 0 \implies A \ginv{A} A \alpha = 0 \implies A\alpha = 0\) and \( A\alpha = 0 \implies \ginv{A} A \alpha = 0 \) . i.e, if a vector \( \alpha \) is orthogonal to the rows of \( A \), then it is orthogonal to the rows of \( H \) and vice versa. Hence \( \rank{A} = \rank{H} \).
\end{proof}

\subsection*{General Solution for a System of Homogenous Equations}
Observe that, \( A(I-H) = A - AH = 0 \), by the first property of g-inverse. i.e, each of the \( n \) columns \( h_{1}, h_{2}, \dots, h_{n} \) of \( I-H \) are orthogonal to the rows of \( A \).

Also, \( (I-H)^{2} = I - H - H + H^{2} = I - H \) by second property of g-inverse. i.e, \( I-H \) is idempotent.\footnote{Let \( R \) be a ring with unity \( 1_{R} \) and \( e \) be an idempotent element of this ring. Then, \( 1_{R} - e \) is also idempotent.}
This implies \( \rank(I-H) = \operatorname{Tr}(I-H) = \operatorname{Tr}(I) - \operatorname{Tr}(H) = n - r \), where \( r = \rank(A) = \rank(H) \) from third property of g-inverse.

Thus, only \( (n-r) \) of the columns vectors \( h_{1}, h_{2}, \dots, h_{n} \) are linearly independent. Without loss of generality we assume that the first \( (n-r) \) columns of \( I-H \), \( h_{1}, h_{2}, \dots, h_{n-r} \) are linearly independent.

Since \( A \) is an \( m \times n \) matrix of rank \( r \), its rows are \( n \)-vectors and therefore we can find at most \( (n-r) \) linearly independent vectors orthogonal to them. \( h_{1}, h_{2}, \dots, h_{n-r} \) is one such set. If there exists any other vector orthogonal to the rows of \( A \), it must be a linear combination of \( h_{1}, h_{2}, \dots, h_{n-r} \).

Let us now consider a system of homogogenous equations
\begin{equation}
	Ax = 0 \label{hom_eq}
\end{equation}

By definition, \( x \) is orthogonal to the rows of \( A \) and so any vector satisfying \eqref{hom_eq} must be a linear combination of \( h_{1}, h_{2}, \dots, h_{n-r} \).
This is equivalent to saying that \( x \) will be a linear combination of columns of \( I-H \), \( h_{1}, h_{2}, \dots, h_{n} \), because \( h_{n-r+1}, h_{n-r+2}, \dots, h_{n} \) are linear combinations of \( h_{1}, h_{2}, \dots, h_{n-r} \) themselves.

Therefore, \( x \) must be of the form
\begin{equation}
	x = z_{1}h_{1} + z_{2}h_{2} + \dots + z_{n}h_{n} = (h_{1}\ h_{2}\ \dots \ h_{n})z = (I-H)z \label{x_lin_eq}
\end{equation}
for some \( z = (z_{1}\ z_{2}\ \dots \ z_{n})' \).\footnote{These \( z_{i} \)'s need not be unique even for a fixed \( x \) as columns of \( I - H \) are not necessarily linearly independent.}

Conversely, if \eqref{x_lin_eq}  holds, \( Ax = A(I-H)z = (A-AH)z = 0 \), by the first property of g-inverse.

This shows that the general solution of \eqref{hom_eq} is given by \( \widetilde{x} = (I-H)z \).

\subsection*{General Solution for a System of Non-Homogenous Equations}
Consider a consistent system of non-homogenous equations given by
\begin{equation}
	Ax = u \label{non_hom_eq}
\end{equation}

If \( \ginv{A} \) is any g-inverse of \( A \), by definition of \( \ginv{A} \), \( \ginv{A}u \) is a particular solution to \eqref{non_hom_eq}. 
Thus \( A(x - \ginv{A}u) = u - u = 0 \) is a system of homogenous equations in \( x - \ginv{A}u \).
The corresponding general solution of this homogenous system is \( \widetilde{x} - \ginv{A}u = (I-H)z \) which implies, the general solution to \eqref{non_hom_eq} is, 
\begin{equation}
	\widetilde{x} = \ginv{A}u + (I-H)z \label{gen_sol}
\end{equation}

\section*{Solution of the Normals Equations (General Form)}
The normal equations are given by (as seen from \eqref{normal_eq}), 
\[ X'Y = (X'X) \htb \] 
A particular solution of these equations will be 
\[ \htb = \ginv{S}X'Y \] where \( \ginv{S} \) is any g-inverse of \( S= X'X\).
Thus, the general solution of normal equations will be, 
\begin{equation}
	\widetilde{\beta} = \htb + (I-H)z \label{gen_sol}
\end{equation}
where \( H = \ginv{S}S \) is a \( p \times p \) matrix.

\subsection*{Some More Results}
\begin{enumerate}[1., start = 3]
	\item \( X = XH \)
\begin{proof}
	From the first property of g-inverse, we have \( SH = S \).Therefore,
	\begin{align*}
		0 &= (I-H)'(S-SH) \\
		  &= (I-H)'(X'X - X'XH) \\
		  &= (I-H)'(X'(X-XH)) \\
		  &= (X(I-H))'(X-XH) \\
		  &= (X-XH)'(X-XH)
	\end{align*}
Note that both the sides of the equality are \( p \times p \) matrices.
Thus equating the \( i \)th diagonal elements of both sides, we get that the sum of squares of the entries in \( i \)th row of \( (X-XH) \) is zero for every \( i = 1, 2, \dots, p \).
Thus every element of \( (X-XH) \) is zero.
\end{proof}

	\item A solution of the normal equations \eqref{normal_eq} is unique iff \( \rank(X) = \rank(X'X) = p \).
\begin{proof}
	A solution to the normal equations is unique iff the general solution \eqref{gen_sol} of normal equations doesn't contain the arbitrary vector.

	This happens iff \( I-H = 0 \) i.e, \( I = H = \ginv{S}S \) which happens iff \( S \) is non-singular and has regular inverse \( S^{-1} \). Thus the result follows.
\end{proof}



\end{enumerate}

\end{document}

